\documentclass[11pt]{article}

% ---- Maintainer: revision authority (outer repo) ----
% Phase 2 prose is AUTHORIZED. Edit ONLY this file for the flagship general-science narrative (no new capstone TeX).
% Normative: SPEC_024_FS1 (specs/.../SPEC_024_FS1_FIELD_SCIENCE_STRATEGIC_CLOSURE.md) §2 field core + §6 integration outline.
% Gates + discipline: docs/007_FORMAL_CLOSURE_MEMO.md. Optional Lean/name evidence: docs/008_ADVISOR_SUBSTANTIVE_CLOSURE_REPORT.md.

% ---- Packages ----
\usepackage{amsmath,amssymb,amsthm}
\usepackage{mathtools}
\usepackage[hidelinks]{hyperref}
\usepackage{cleveref}
\usepackage{booktabs}
\usepackage{enumitem}
\usepackage[margin=1.15in]{geometry}
\usepackage{microtype}
\usepackage{xcolor}
\usepackage{tcolorbox}

% ---- Key-point boxes ----
\newtcolorbox{keybox}[1][]{
  colback=blue!4, colframe=blue!35,
  arc=3pt, boxrule=0.5pt,
  fonttitle=\bfseries, #1
}

% ---- Lean name formatting (line-breakable) ----
\ExplSyntaxOn
\tl_new:N \l__rfo_leanref_tl
\NewDocumentCommand \leanref { m }
  {
    \tl_set:Nn \l__rfo_leanref_tl { #1 }
    \regex_replace_all:nnN { \c{_} } { \c{_} \c{allowbreak} } \l__rfo_leanref_tl
    \regex_replace_all:nnN { \. } { . \c{allowbreak} } \l__rfo_leanref_tl
    \regex_replace_all:nnN { ([a-z])([A-Z]) } { \1 \c{allowbreak} \2 } \l__rfo_leanref_tl
    \regex_replace_all:nnN { ([0-9])([A-Z]) } { \1 \c{allowbreak} \2 } \l__rfo_leanref_tl
    \texttt{\small \tl_use:N \l__rfo_leanref_tl }
  }
\ExplSyntaxOff

% ---- Theorem environments ----
\theoremstyle{plain}
\newtheorem{theorem}{Theorem}[section]
\newtheorem{lemma}[theorem]{Lemma}
\newtheorem{corollary}[theorem]{Corollary}
\newtheorem{proposition}[theorem]{Proposition}

\theoremstyle{definition}
\newtheorem{definition}[theorem]{Definition}
\newtheorem{example}[theorem]{Example}

\theoremstyle{remark}
\newtheorem{remark}[theorem]{Remark}

% ---- Macros ----
\newcommand{\RFO}{\textsc{rfo}}
\newcommand{\RI}{\textsc{ri}}
\newcommand{\Lean}{\textsc{Lean\,4}}
\newcommand{\Mathlib}{\textsc{Mathlib}}
\newcommand{\ReflTransGen}{\operatorname{ReflTransGen}}
\newcommand{\PreservedBy}{\operatorname{PreservedBy}}
\newcommand{\ForwardClosed}{\operatorname{ForwardClosed}}
\newcommand{\reachableFrom}{\operatorname{reachableFrom}}
\newcommand{\unionGen}{\operatorname{unionGen}}
\newcommand{\obj}{\mathsf{obj}}
\newcommand{\mor}{\mathsf{mor}}
\newcommand{\represent}{\mathsf{represent}}
\newcommand{\ProvBic}{\mathrm{ProvBic}}

\title{The General Science of Reflexive Systems:\\
Nonexhaustibility, Completion Failure, and Residual Structure}
\author{Nova Spivack\\Independent Research}
\date{}

\begin{document}
\maketitle

\begin{abstract}
This paper develops a general science of reflexive systems: systems that, in a mathematically nontrivial sense, represent, evaluate, certify, locate, or attempt to complete themselves from within. The same structural tensions recur wherever serious self-representation, internal generation, and self-certification meet. Its central claim is that strong \emph{anchored} internal completion is structurally limited. More precisely, when a reflexive system attempts anchored internal completion under sufficiently strong representational, generative, and semantic demands, completion is blocked by lawful barrier families. The failure is not mere collapse. It leaves a structured residual aftermath, with typed witness forms, enrichment channels, and semantic barrier regimes.

The paper organizes this science around six interlocking components: completion modes and anchored vs relocated success; flagship anchored impossibility under three barriers; forced aftermath with witnesses and enrichment; two semantic summits (predicated computational vs proof-theoretic full); explicit alignment of summit strength back into the residual pipeline; and a \emph{dynamical} layer---the Reflexive Development Law---showing that standing residual is not inert but drives lawful response: refinement, regime shift, or fold-obstruction transition.

The outcome is a single scientific arc---not a catalogue of unrelated limitative tricks. Full anchored self-exhaustion in the sense ruled out here would collapse reflexive depth; the paper's positive payload is what lawfully replaces that fantasy, including an account of how reflexive systems \emph{continue} under barrier discipline rather than merely halting at impossibility.
\end{abstract}

\begin{keybox}[title=Field law (one sentence)]
Reflexive systems are governed not by perfect internal completion, but by the lawful interplay of completion ambition, barriered failure, structured residuality, and dynamical response under standing residual burden.
\end{keybox}

\section{Introduction}

\subsection{A new scientific subject}

There is a broad class of systems whose behavior cannot be understood solely by describing external inputs, outputs, and transition laws. These are systems that stand in a nontrivial internal relation to themselves. They represent themselves, evaluate their own states, certify or invalidate their own claims, construct internal surrogates of their own operation, or attempt some form of internal completion. Formal systems with truth and proof predicates, universal computational systems, self-modeling agents, reflective semantic architectures, and reflexive epistemic networks all belong, in different ways, to this family.

This paper studies that family under a single heading:

\begin{definition}[Reflexive system]
A \emph{reflexive system} is a system whose operation includes a mathematically nontrivial internal relation to itself: self-representation, self-evaluation, self-certification, self-location, or attempted internal completion.
\end{definition}

The central thesis of the paper is that reflexive systems are governed by a common structural law. They can develop rich internal relations to themselves, but sufficiently strong internal completion is blocked. This blockage is not a mere pathology of one theorem or one formalism. It is a domain law. Moreover, the blockage is not the end of the story: reflexive systems that continue under barrier discipline must \emph{respond} to the standing residual---through refinement, architectural reconfiguration, or regime transition---and this response is itself lawful.

\begin{keybox}
The central scientific claim of this paper is not simply that reflexive systems can fail at self-completion. It is that their failure is lawful, classifiable, and residually structured.
\end{keybox}

\subsection{Why reflexive systems form a natural scientific class}

A skeptical reader may fear that ``reflexive system'' is a label invented to package one research program. The reply is structural. Across logic, computability, semantics, and engineering, one meets systems whose \emph{internal} relation to themselves is part of the explanatory subject, not an optional gloss:
\begin{itemize}[leftmargin=1.5em]
\item \textbf{Formal truth and proof.} Theories that encode their own syntax and provability, or that quantify over codes for their own assertions, exhibit reflexive organization essentially (Gödel, Tarski tradition).
\item \textbf{Computability and fixed points.} Universal machines, self-interpreters, and Kleene-style recursion theorems produce nontrivial self-reference as a \emph{mathematical mechanism}, not an accident of notation.
\item \textbf{Semantic closure and certification.} Architectures that judge, score, or certify their own outputs instantiate reflexive evaluation---from truth predicates (when available) to program logics and verification pipelines.
\item \textbf{Self-modeling agents.} Any system that maintains an internal model indexed to its own states or policies stands in a reflexive relation that can be made mathematically precise.
\item \textbf{Epistemic and institutional reflexivity.} Scientific and organizational systems that represent their own methods, standards, or error modes repeat the same pattern at a different scale.
\end{itemize}

What unifies these cases is not a shared implementation stack but a shared \emph{problem form}: how can internal representation, internal generation, and internal evaluation cohere without collapsing into trivial self-coincidence or false closure? That problem form appears independently in many domains. A science that names it, classifies completion attempts, and proves lawful failure-and-residue patterns is therefore responding to a recurring natural structure, not merely narrating a private formal artifact.

\subsection{From incompleteness results to a science of reflexive systems}

The classical limits due to Gödel, Tarski, Kleene, Lawvere, and their successors are usually stated as \emph{negative} theorems: no comprehensive self-certain truth definition, no finitely axiomatized self-representing theory that proves its own consistency in the naive sense, no unrestricted diagonal-free totality of computable self-transformations, no degenerate categorical fixed-point collapse without cost.

This paper does \emph{not} replace those traditions. It reorganizes their moral. Gödel exposed the limits of ambitious formal self-totalization; Tarski exposed semantic non-self-closure; Kleene exposed how fixed-point and recursion-theoretic internality constrain what can be uniformly realized; Lawvere exposed how diagonal structure permeates categorical settings. Together, they show that reflexive organization is \emph{costly}.

What has been missing at the level of \emph{science}---as opposed to a list of limitative results---is a unified account of \emph{(i)} completion attempts, \emph{(ii)} lawful anchored barriers, \emph{(iii)} forced positive aftermath, \emph{(iv)} semantic stratification of closure strength, and \emph{(v)} lawful dynamical response to standing residual. That quintet is the positive scientific payload; \S13 collects its canonical statement once the arc is in view. The flagship impossibility theorem opens the arc; what follows---including the dynamical layer---is not optional ornament but the \emph{same} science continuing under lawful failure.

\subsection{From impossibility to science}

Much of the classical literature on self-reference identifies impossibility phenomena: undefinability, incompleteness, diagonal barriers, fixed-point pathologies, undecidable fragments, and semantic escape. These results are profound, but they often stop at the negative limit. They tell us that certain closure ambitions fail, but not what the ambient science of such failure should be.

Here the aim is broader. We seek a theory of:
\begin{enumerate}[label=(\arabic*)]
\item what kinds of internal completion reflexive systems can attempt,
\item why those attempts fail under explicit barrier interfaces,
\item what structured aftermath remains when they fail,
\item how semantic barrier regimes stratify the possibilities,
\item how these regimes integrate with the residual theory,
\item and how reflexive systems \emph{respond} to standing residual---through refinement, regime shift, or fold---under lawful dynamical constraints.
\end{enumerate}

This shift matters. A mature science of reflexive systems should not only identify impossibility. It should identify the lawful aftermath of impossibility and the lawful dynamics that follow.

\subsection{Main contributions}

The paper makes six central contributions.

\begin{enumerate}[leftmargin=2em]
\item It develops a general classification of completion ambitions in reflexive systems, distinguishing representational, generative/closure, and semantic/certificatory completion modes.
\item It proves a flagship theorem: under the three obstruction interfaces, true honest anchored internal completion is impossible.
\item It proves that the failure of completion has lawful aftermath, yielding a positive residual profile and barrier-linked residual witnesses.
\item It constructs enriched residual witnesses carrying representational, closure, certification, and mixed obstruction payloads.
\item It develops two semantic barrier summits --- a computational predicated summit and a stronger proof-theoretic summit --- and integrates them into the residual theory.
\item It establishes a dynamical layer: the \emph{Reflexive Development Law}, showing that standing residual drives lawful response---refinement, proper regime shift, or fold-obstruction transition---and that adequacy burden relocates rather than vanishes under honest dynamical episodes.
\end{enumerate}

\subsection{What this paper is not}

This is not a repository tour or a report on a formalization project. The formal artifacts matter because they sharpen definitions, block smuggling, and force the theory to become explicit. But the paper's subject is broader: it proposes a general science of reflexive systems, using formal realization to establish its claims with unusual precision. Readers outside the formal stack should treat the machine-checked layer as an audit trail, not as the \emph{source} of the scientific problems.

\section{Reflexive systems and internal completion}

\subsection{Completion as the governing ambition}

The central organizing notion of the paper is not merely self-reference but \emph{internal completion}. Reflexive systems differ in many ways, but the most important divide is whether they attempt to complete themselves from within. Some systems merely mirror themselves; others attempt to exhaust the relevant internal burden.

\begin{definition}[Internal completion]
An \emph{internal completion attempt} is an internal attempt by a reflexive system to exhaust, determine, certify, generate, or semantically capture the relevant self-related burden from within.
\end{definition}

This notion is intentionally broad at the conceptual level. The scientific work begins when we ask how such completion can be attempted.

\subsection{Three principal completion modes}

Across reflexive practice, internal completion pressure organizes into three \emph{central structural families}. These are not mere bookkeeping columns invented for this paper; they are the principal ways serious systems attempt to settle self-related burden internally.

\begin{definition}[Representational completion]
A system attempts \emph{representational completion} when it tries to capture itself through an internal representational surrogate, image, code, or model that is supposed to exhaust the relevant burden.
\end{definition}

\begin{definition}[Closure completion]
A system attempts \emph{closure completion} when it tries to generate or close all relevant internal transitions, reaches, folds, or continuation structure through an internal generative or closure mechanism.
\end{definition}

\begin{definition}[Semantic/certificatory completion]
A system attempts \emph{semantic or certificatory completion} when it tries to decide, certify, or semantically totalize the relevant self-related content from within.
\end{definition}

These three families are not arbitrary. They correspond to the three great ways a reflexive system can claim that its self-related burden is settled:
\begin{itemize}
\item by being internally represented,
\item by being internally generated or closed,
\item or by being internally certified or semantically exhausted.
\end{itemize}

\subsection{Anchored completion versus relocated success}

A major conceptual distinction is required here.

\begin{definition}[Anchored completion]
An internal completion attempt is \emph{anchored} when the relevant success is borne at the internal point, locus, or anchor at which completion is supposed to occur.
\end{definition}

\begin{definition}[Relocated success]
A completion-like phenomenon is \emph{relocated} when success occurs only away from the relevant anchor, or only somewhere else in a larger parameter space, carrier, or auxiliary geometry.
\end{definition}

This distinction is not terminological. It changes the theorem subject. If success is allowed to relocate, then many impossibility theorems for anchored completion do not apply. But this does not refute those theorems. It changes the object under study.

\begin{keybox}
Anchored completion and relocated success are different scientific subjects. A theorem about anchored impossibility is not defeated by a weaker off-anchor success phenomenon.
\end{keybox}

The science of reflexive systems must therefore distinguish at least two layers:
\begin{enumerate}[label=(\roman*)]
\item strong anchored completion,
\item weaker geometries of relocated or off-anchor success.
\end{enumerate}

The flagship theorems of this paper concern the first layer.

\section{The barriered impossibility of anchored completion}

\subsection{Why anchored completion is the central completion notion}

The flagship theorem is stated for \emph{anchored} honest internal completion, not for every conceivable success predicate. That choice is scientifically central.

\paragraph{Strongest serious notion.}
Among internal completion attempts, anchored completion is the notion that disciplines representation, generation, and evaluation at the locus where completion is \emph{supposed} to discharge the burden. It is the modality that corresponds to ``the system actually settles its own reflexive debt at the appointed place,'' not ``success can be found if we widen the geometry.'' If one were to prove impossibility only for deliberately weakened or relocation-dependent targets, skeptical readers would rightly ask whether anything substantive had been shown.

\paragraph{Relocated success is a different subject.}
When success is allowed to migrate---to occur only off a designated anchor, or only in an enlarged carrier---the theorem subject changes. Many classical phenomena are best read as \emph{relocated} rather than \emph{anchored} success geometries. The paper does not treat relocation as a loophole that refutes anchored impossibility; it treats it as a distinct scientific layer whose classification belongs in the same science but under different hypotheses.

\paragraph{Why the three barrier families are scientifically central.}
The three interfaces correspond to the three dominant ways a reflexive system might claim its self-related burden is settled: internal representation, internal generation/closure, internal certification or semantic exhaustion. The paper does \emph{not} claim a single formal theorem that enumerates \emph{every} imaginable completion gadget in all possible universes. What it does claim is that, for scientific purposes, organizing obstruction along these three families captures the principal axes along which serious completion attempts have been understood from Hilbert-era optimism through modern reflexive semantics and computer-assisted certification. Proving anchored impossibility under simultaneous obstruction along all three axes therefore addresses the conjunction of the main classical pressures, not a narrow ad hoc trio.

\subsection{Three obstruction interfaces}

The nonexistence theorem is organized around three obstruction families.

\begin{definition}[Representational obstruction]
A \emph{representational obstruction interface} expresses the impossibility of successful internal representational completion.
\end{definition}

\begin{definition}[Closure obstruction]
A \emph{closure obstruction interface} expresses the impossibility of successful internal generative or closure completion.
\end{definition}

\begin{definition}[Semantic obstruction]
A \emph{semantic obstruction interface} expresses the impossibility of successful internal semantic or certificatory completion.
\end{definition}

The abstract formal stack packages these as the three barrier interfaces on a common reflexive architecture. Their scientific meaning is simple: each says that a corresponding mode of anchored completion cannot succeed.

\subsection{Honest anchored completion}

The target of the flagship theorem is not arbitrary self-related data but \emph{honest anchored internal completion}. In the formal realization this is represented by an anchored completion package with a true internal completion claim plus a burden-bearing discipline. The important point is conceptual: we are not ruling out every object carrying reflexive structure. We are ruling out the existence of a true burden-bearing anchored completion package under the three barriers.

\subsection{Flagship theorem}

\begin{theorem}[Barriered anchored completion is impossible]\label{theorem:barriered-anchored}
Let \(A\) be a reflexive architecture satisfying the representational, closure, and semantic obstruction interfaces. Then there does not exist a true honest anchored internal completion package over \(A\).
\end{theorem}

\begin{proof}[Proof architecture]
The proof has four steps.

First, a true honest anchored completion package forces anchored completion-mode structure. It cannot remain semantically inert; it must commit to one of the canonical anchored completion families.

Second, each anchored family induces a corresponding success predicate at the architecture level: representational success, closure success, or semantic/certificatory success.

Third, the three obstruction interfaces jointly forbid all such success predicates.

Fourth, contradiction follows. Hence no true honest anchored completion package exists.
\end{proof}

This is the correct theorem surface for the paper. The four-way cover and related profile decompositions are part of the internal proof architecture, but the scientific theorem is global nonexistence.

\begin{keybox}
The public theorem is not ``one of four disjuncts holds.'' It is: under the three barriers, true honest anchored completion cannot exist.
\end{keybox}

\subsection{Interpretation: which reflexive aspirations are ruled out?}

Theorem~\ref{theorem:barriered-anchored} (informally, barriered anchored completion is impossible) does not assert that reflexive systems cannot \emph{do} anything interesting. It rules out a specific joint aspiration: \emph{simultaneous} honest settlement of representational, generative/closure, and semantic/certificatory burden \emph{at the anchor} in a way that would amount to full anchored internal completion under the stated interfaces.

What is ruled out includes idealized templates such as: a reflexive architecture that internally represents itself faithfully enough to dissolve diagonal tension, closes all relevant internal generation at the anchor, and certifies or semantically exhausts its self-related claims---all while carrying a true, burden-faithful completion package. What is \emph{not} ruled out, without further argument, is weaker or relocated success, partial satisfiability of isolated barriers, or rich dynamics in the residual regime once the conjunction fails.

\subsection{Why this is stronger than a local diagonal no-go}

The theorem should not be read as merely another diagonal paradox. It is a law about a broad class of burden-bearing internal completion attempts. It says that if a reflexive system tries to settle its own self-related burden in the strong anchored sense, then under the three obstruction interfaces that project fails globally.

\paragraph{Why the science cannot stop here.}
A science of reflexive systems that ended at impossibility would mistake the subject for a catalogue of gloomy limits. In domains that study living, adaptive, or semantically rich reflexivity, \emph{what happens next} is as important as the failure itself. If strong anchored completion is blocked, the honest question is not only ``blocked how?'' but ``what lawful state succeeds the block, and what happens when a reflexive system \emph{continues} under that block?'' Answering those questions is not a digression from the flagship theorem; it is the inevitable continuation of the same inquiry---through aftermath, witness enrichment, semantic stratification, and ultimately through the \emph{dynamical} response layer developed later in this paper.

\section{Aftermath: residual structure is forced}

\noindent\textbf{Plain language.} \textbf{Residual structure} is the lawfully organized remainder of a reflexive system's self-related burden after strong internal completion has failed. It is not arbitrary debris: it is where reflexive organization continues under explicit constraints once closure ambitions break.

\noindent\emph{Intuition.} Equivalently: residual structure is the part of that burden that remains \emph{typed, witness-linked, and integrable} into richer obstruction payloads---closer to a \emph{live} unresolved configuration than to a dead remainder.

\subsection{Failure is not collapse}

The next major claim is positive.

\begin{theorem}[Structured aftermath]
If the canonical anchored completion modes are blocked, then the resulting state is not merely null or chaotic. A positive residual profile is forced.
\end{theorem}

The point is that completion failure does not end the theory. It begins the next layer of the science.

\begin{keybox}[title=Aftermath principle]
Failure is not a defect of reflexive systems in the sense relevant here. It is one of the conditions under which nontrivial reflexive organization can remain \emph{deep} rather than collapsing into pretend closure.
\end{keybox}

\subsection{Residual profile}

The formal stack isolates a positive residual profile that appears precisely after the failure of the canonical completion modes: the reflexive system is not internally completed, but its noncompletion is itself organized.

\subsection{Why residual structure is not a fourth completion mode}

\begin{remark}
The residual branch is not a fourth completion mode. The first three families are completion ambitions; the residual branch is an aftermath family appearing once those ambitions fail. Blurring this distinction would collapse the science into confusion.
\end{remark}

\subsection{Residual witnesses}

The aftermath is not merely asserted propositionally. It is witnessed. In the formal layer, this appears as \(R_4\)-class residual witnesses that are linked to the barrier pack.

Scientifically, the significance is that failure carries typed residual data rather than merely negative absence.

\section{Why nonexhaustibility is not a defect}

The flagship impossibility theorem can be misread as a gloomy conclusion: ``reflexive systems cannot win.'' The correct reading is almost the opposite.

\subsection{Total completion would trivialize reflexive depth}

A reflexive system that were \emph{fully} internally exhausted, in the strong anchored sense ruled out by the barrier interfaces, would no longer carry the tension that makes reflexive structure scientifically interesting. Total self-settlement at the anchor, across representation, generation, and certification, would amount to a false paradise: a self that has, in effect, nothing left to reflexively negotiate. The impossibility theorem says that such paradise is not available as an honest anchored package under the stated laws---not that reflexive life is impoverished, but that it cannot be frozen into a self-transparent totality.

\subsection{Failure preserves nontriviality}

Barriered failure is what prevents collapse into degenerate self-coincidence. Without lawful obstruction, the only alternative is often smuggling: smuggled assumptions, hidden parameters, or informal moves that pretend closure while evading explicit burden. The barriers force honesty about where completion succeeds or fails. In that sense, obstruction is a \emph{guardrail} for nontrivial reflexive science.

\subsection{Kinds of failure: bug, boundary, constitution}

It helps to separate three rhetorical senses of ``failure'':
\begin{itemize}[leftmargin=1.5em]
\item \textbf{Bug-like failure} is pathological: an implementation error, an unsound rule, or an accidental inconsistency. Removing the bug restores intended behavior.
\item \textbf{Boundary-like failure} marks the edge of a regime: a statement is independent, a resource bound bites, a class restriction is necessary. The system still behaves as designed; the lesson is cartography.
\item \textbf{Constitutive failure} is the sense central to this paper. Here failure is not a flaw to eliminate but a \emph{lawful outcome} of serious reflexive organization under explicit interfaces. The residual aftermath is not debris to sweep away; it is where structured reflexive life continues after completion ambition meets barrier.
\end{itemize}

The thesis of this section is that reflexive nonexhaustibility, in the sense developed here, is closest to \textbf{constitutive} failure. The theory's positive energy lies in the aftermath and enrichment layers, not in hoping the barriers disappear.

\subsection{Residual structure as ongoing reflexive life}

Residual profiles and enriched witnesses are not consolation prizes. They are the mathematical forms of \emph{continued} reflexive structure: types of obstruction that survive, can be witnessed, and can be enriched. Nonexhaustibility thus names not a wound but the \emph{space in which} a reflexive system can keep doing honest reflexive work. The next question is unavoidable: \emph{how} is that organized aftermath carried?

\section{Residual witnesses and enrichment}

Witness theory turns aftermath from a blanket existence claim into operational content: \emph{how} is the forced profile materialized as typed data, and how does obstruction survive in distinct columns?

\subsection{Barrier-linked residual witnesses}

The first witness layer is abstract. Under the barrier interfaces, an \(R_4\)-witness can be constructed whose carrier literally contains the barrier pack. This is the abstract D-002 level.

\subsection{Single-column enrichment}

The abstract witness can be enriched along three distinct obstruction columns:
\begin{enumerate}[label=(\arabic*)]
\item a representational obstruction payload,
\item a closure obstruction payload,
\item a certification obstruction payload.
\end{enumerate}

These are not decorative tags. Each column carries typed obstruction content.

\subsection{Mixed enrichment}

When the three barriers hold simultaneously, the witness can be enriched with all three columns at once in a named three-field mixed payload. This yields a stronger post-failure theorem:

\begin{theorem}[Mixed residual enrichment]
Under simultaneous representational, closure, and semantic obstruction, the residual witness carries mixed obstruction content from all three columns at once.
\end{theorem}

This is a major strengthening of the aftermath theory. It shows that post-failure structure is not merely nonempty but internally differentiated.

\subsection{Interpretation: what to expect architecturally}

If anchored internal completion is impossible under the joint barriers but structured residuality is lawful, then architectures should be designed for \emph{explicit} barrier discipline rather than for fantasy total self-settlement. Expect:
\begin{itemize}[leftmargin=1.5em]
\item typed witness layers that record how each obstruction column fails or persists,
\item separation between anchored obligations and relocated or auxiliary success,
\item enrichment pipelines that preserve obstruction content rather than masking it,
\item and semantic stacks that declare whether they live at predicated or full barrier strength.
\end{itemize}
The science discourages ``paper over the tension'' engineering; it rewards reflexive honesty.

\section{Semantic barrier regimes I: the computational predicated summit}

Witnesses describe \emph{how} residue is carried. They do not yet specify \emph{how strong} self-semantic barrier closure can be in different organizational regimes. Without summit stratification, one cannot distinguish honest class-restricted strength from delusions of unrestricted universality.

\subsection{The two principal summit regimes (scope)}

This paper develops two fully articulated \emph{semantic barrier} summit regimes: a \textbf{computational predicated} summit and a \textbf{proof-theoretic full} summit. They are the first two such regimes worked out to closure in the present science with machine-checked summit interfaces. The paper does \textbf{not} claim that no further semantic regimes will ever matter---only that these two are structurally central \emph{now}: they crystallize the distinction between class-restricted (predicated) closure and stronger proof-theoretic full closure, and they supply the comparison and integration story that follows.

Readers should therefore treat them as \textbf{canonical exemplars} and as \textbf{scientifically obligatory} distinctions within reflexive barrier theory, without interpreting the text as asserting metaphysical exhaustiveness over all future semantics.

\subsection{A concrete computational frame}

The first route is computational. A concrete self-semantic frame is built on a Kleene/\(\texttt{Nat.Partrec.Code}\)-style substrate. The crucial features are:
\begin{itemize}
\item concrete code,
\item a nontrivial extensional equivalence on code,
\item invariant realized truth,
\item encoded nontriviality,
\item and a class-restricted semantic barrier package.
\end{itemize}

\subsection{Predicated barrier hypotheses}

The outcome is not full unrestricted barrier closure but a predicated summit.

\begin{theorem}[Computational predicated summit]
The Kleene self-semantic frame yields a nontrivial self-semantic system satisfying a predicated semantic barrier package for the computable transformer class.
\end{theorem}

This is the correct shape of the computational route. The summit is \textbf{not} weak because it is ``sloppy'' or half-finished. On the contrary: for the \textbf{computable transformer class} it governs, it achieves real, disciplined barrier closure---the restriction to that class is the \emph{principled} content of the route, not an accident to overcome. The obstruction analysis shows exactly why one must not expect, from that route alone, to obtain the same unrestricted universal closure demanded by full \(\mathrm{BarrierHypotheses}\) for all code transformers.

\subsection{What this teaches}

The computational route is not a failed attempt at the stronger summit. It is a distinct and lawful summit whose scientific role is \textbf{restricted strength}: strong on a class, honest about the boundary. \textbf{Positive reading:} the Kleene summit is a \emph{genuine} semantic achievement on the computable transformer class---a real closure regime, not an approximate hack awaiting ``the real theorem.'' Its scientific office is to mark exactly where that achievement ends and why further universalization demands a different semantic organization. The hierarchy with the proof-theoretic route is therefore not ``bad versus better'' in a moralized sense, but \textbf{scientifically stratified strength}:
\[
\text{computational restricted closure} \neq \text{proof-theoretic full closure}.
\]

\section{Semantic barrier regimes II: the proof-theoretic full summit}

\subsection{Need for a stronger route}

To reach full semantic barrier closure, the computational line is not enough. A stronger route is required.

\subsection{The proof-theoretic frame}

The proof-theoretic route is built from:
\begin{itemize}
\item a Gödel system,
\item a middle equivalence of the \(\ProvBic\)-type,
\item arithmetical semantics,
\item encoded nontriviality,
\item representability and reflection structure.
\end{itemize}

\subsection{Main theorem}

\begin{theorem}[Proof-theoretic semantic barrier summit]
Gödel-system structure together with \(\ProvBic\)-arithmetical semantics yields full \(\mathrm{BarrierHypotheses}\).
\end{theorem}

This theorem is stronger than the computational predicated summit and resolves the earlier obstruction in the correct way: not by claiming that the bare Gödel shell implies full barrier closure, but by adding the semantic layer that makes the full package meaningful.

\textit{Scientific significance.} Full semantic barrier closure is \textbf{not vacuous in principle}: it is \textbf{attainable} under proof-theoretic organization strong enough to support the full package, even though it is unavailable on the current computational API path. The computational route's limitation is structural, not a moral that ``nothing ever closes.'' The contrast locates where genuine universal-grade barrier strength honestly lives.

\begin{keybox}[title=Depth consequence]
A reflexive system that were fully internally exhausted in the strong anchored sense ruled out here would cease to be reflexively deep in the scientifically relevant sense. Nonexhaustibility is not a defect of reflexive systems but one of the conditions of their depth.
\end{keybox}

\subsection{Interpretation}

The proof-theoretic route shows that full semantic barrier closure is achievable in a disciplined form, but only at a stronger semantic level than the purely computational summit.

\section{Comparing the semantic summits}

\subsection{Comparison at the API layer}

The computational and proof-theoretic summits should not be conflated. They are not isomorphic as self-semantic frames. The comparison belongs at the barrier-hypothesis/API layer.

\subsection{The computational summit is genuinely weaker}

The formal comparison proves that the Kleene computational route does not promote to the full summit on its present API path. This is not a missing theorem; it is a structural fact.

\begin{theorem}[Computational summit does not upgrade for free]
The current Kleene predicated summit does not, on its own API path, upgrade to full \(\mathrm{BarrierHypotheses}\).
\end{theorem}

\subsection{The proof-theoretic summit has a predicated shadow}

Conversely, the full proof-theoretic summit has a predicated shadow. So the relation is asymmetric in the correct direction:
\[
\text{full summit} \Longrightarrow \text{predicated shadow},
\qquad
\text{but not conversely on the current computational route}.
\]

\subsection{Scientific significance}

This comparison stabilizes the semantic barrier science. The two summits are related, but not by collapse. The lesson is not merely that one route ranks below another on a single axis. It is that \textbf{semantic barrier strength is itself a stratified resource} of reflexive systems: different organizations yield different honest ceilings on closure.

\subsection{Implications for self-certifying and reflective systems}

The stratification between predicated and full semantic barrier strength has direct readings for systems that verify, certify, or reflect on themselves:
\begin{itemize}[leftmargin=1.5em]
\item \textbf{Self-certifying pipelines.} Many real stacks implicitly live at a predicated or resource-aware strength: they prove properties for a restricted class of updates or transformations. The theory says that upgrading to full strength is not notation-deep; it may require a different semantic regime (here, the proof-theoretic route).
\item \textbf{Reflective AI and semantic verification architectures.} When a model or service evaluates its own outputs or maintains an internal consistency policy, the predicated/full distinction warns against conflating ``works on the guarded class we tested'' with ``full reflexive closure.''
\item \textbf{Formal scientific systems.} Libraries of mathematics and proof assistants already embody barrier discipline; the paper's science explains why residual structure and enrichment channels are predictable features, not accidents, when completion ambition is honest.
\end{itemize}

\paragraph{Why integration is not the final word.}
Integration shows that impossibility, semantics, and residue are \textbf{one system}---not three bibliographies forced between the same covers. But the science still owes an answer: \emph{what happens next?} That is the dynamical layer developed after the integration arc.

\section{Integration into the residual theory}

\subsection{Why integration is necessary}

Semantic barrier summits would remain scientifically \emph{isolated} without a path back into the theory of what survives completion failure; residual theory would be \emph{semantically blind} without articulating how different barrier strengths participate in the same enriched aftermath pipeline. Integration is what keeps impossibility, semantics, and residue from fracturing into unrelated specialties.

\subsection{Integration seams (O7)}

The following constructions concretely realize the unity described above: proof-theoretic barrier strength is carried into the same enriched \(R_4\) pipeline that packages post-completion obstruction. The triple-augment seam below is the current capstone: the full paper-side barrier pack meets the summit line inside the same witness pipeline that governs aftermath.

\subsection{O7 constant seam}

The first integration seam is constant in the barrier pack. It is the minimal honest bridge from the proof-theoretic summit into the enriched \(R_4\) pipeline.

\subsection{Repr-dependent seam}

The next seam depends genuinely on the representational barrier field \(b.\mathrm{reprBarrier}\). This proves that dependence on paper-barrier content can enter the semantic barrier package nontrivially.

\subsection{Full triple-augment seam}

The strongest current integration theorem augments the proof-theoretic barrier package by all three \(U123\) interface propositions:

\begin{theorem}[Triple-augment proof-theoretic O7 integration]
Given the full \(U123\) barrier pack, the proof-theoretic barrier package can be augmented by representational, closure, and certification obstruction content and pushed through the enriched \(R_4\) witness pipeline.
\end{theorem}

The conceptual point is delicate and important. The augmenting conjuncts are abstract interface propositions. This does \emph{not} mean that the internal proof-theoretic semantics has magically become native to closure sets or certification functions. It means that the barrier package carried into the integration seam is strengthened in a disciplined way by the full paper-side barrier pack.

\begin{keybox}
The O7 integration arc shows that semantic barriers are not isolated summit objects. They can be fed into the residual architecture of reflexive systems in increasingly strong ways: constant, repr-dependent, and fully triple-augmented.
\end{keybox}

\section{Dynamics: the Reflexive Development Law}

The science so far is \emph{static}: it classifies completion modes, proves anchored impossibility, organizes aftermath, and stratifies semantic regimes. A general science of reflexive systems must also say what happens when a system \emph{continues} under barrier discipline after forced residual formation. This section develops that \emph{dynamical} layer. The present dynamical theory is intentionally bounded: it establishes a lawful response architecture under standing residual burden, not a universal theorem that every residual in every reflexive system is generative in the strongest imaginable sense.

The dynamic layer proves three things: under standing residual burden, reflexive response is structurally constrained to refinement, proper regime shift, or bookkeeping reconfiguration; where iterate-backed closure fails, fold obstruction forces closure obstruction; and adequacy burden relocates rather than vanishes under honest dynamical episodes.

\subsection{Residual persistence}

Residual is not transient noise but a \emph{standing} unresolved burden.

\begin{definition}[Standing residual burden]
A proposition \(P\) names a \emph{standing residual burden} when it records a non-trivially organized aftermath configuration that persists as an obligation for the system under honest barrier discipline.
\end{definition}

\noindent Kernel witnesses---fiber collisions from forgetful maps, infinity-compression geometry, or paired obstruction data---persist as nontrivial standing burden so long as the relevant identification is not fully separated by a refinement step. Non-resolving refinement leaves standing burden intact.

\subsection{The response trilemma (Reflexive Development Law)}

When a reflexive system responds to standing residual, the response must take one of three lawful forms.

\begin{theorem}[Reflexive Development Law --- standing trilemma]\label{theorem:rdl-trilemma}
Let a reflexive system undergo a response episode (formalized as a \emph{residual response step}). Then, classically, the episode is one of:
\begin{enumerate}[label=(\roman*)]
\item \textbf{Refinement:} the system refines its observational identification (finer distinctions; same carrier), and a standing residual burden of the form \(K \lor P \lor B\) persists;
\item \textbf{Proper regime shift:} the architecture itself changes nontrivially (distinct regulatory parameters, not merely relabeled);
\item \textbf{Bookkeeping reconfiguration:} the architecture is formally replaced but with the same underlying structure.
\end{enumerate}
These exhaust the structural alternatives.
\end{theorem}

\noindent The formal realization is \leanref{reflexive\_development\_law\_standing\_trilemma}. There is no fourth option in which standing residual simply disappears without trace.

\paragraph{Content of the refinement branch.}
When standing residual arises from a forgetful kernel, the refinement branch carries an admissible \(R_2\) residual witness and, when barrier data is available, a joint \(R_2\)/\(R_4\) coexistence certificate.

\subsection{Fold obstruction and regime transition}

The second dynamical principle addresses what happens when residual \emph{cannot} be absorbed within the current closure regime.

\begin{theorem}[Fold obstruction forces closure obstruction]\label{theorem:fold-obstruction}
If a reflexive architecture commits to iterate-backed closure adequacy---meaning every closure success is witnessed by internal iterate reachability from a seed \(S_0\) to a target \(B\)---then whenever \emph{no} closure operator can reach \(B\) from \(S_0\) by iteration, the architecture satisfies a closure obstruction interface.
\end{theorem}

\noindent The formal realization is \leanref{reflexive\_development\_law\_fold\_obstruction\_iterate\_geometry}. Iterate-backed adequacy is a genuine \emph{commitment}, not a free property: there exist architectures where every closure operator is declared successful but no iterate-backed adequacy holds for any seed/target pair (\leanref{exists\_Reflexive\-Architecture\_forall\_closureSuccess\_not\_closureIterate\-Adequacy}).

\subsection{Adequacy burden relocation}

Under the triple barriers, the aftermath theorem (\S4) forces admissible \(R_4\) witnesses. The adequacy bridge says that this forced burden \emph{relocates} under architecture cast and regime alignment, rather than vanishing:

\begin{theorem}[Adequacy burden relocation]
Under \(U_{123}\) barrier data on an opaque totalization attempt, the forced admissible \(R_4\) aftermath constitutes standing residual burden. This burden transports across architecture casts and regime-snapshot alignment.
\end{theorem}

\noindent The formal realizations are \leanref{d3\_adequacy\_aftermath\_standing\_burden} and its entry-point family. Weakened \(R_4\)-core variants confirm that burden persists even when only the existential ``there is an admissible \(R_4\) witness'' is retained.

\subsection{Anti-smuggling: what the dynamical layer does not claim}

\begin{itemize}[leftmargin=1.5em]
\item It does \textbf{not} assert that every residual is generative or that every residual forces nontrivial refinement.
\item It does \textbf{not} identify ``self-improvement'' with ``residual response.''
\item It does \textbf{not} supply a universal quantifier ``for all reflexive architectures, the development law holds without hypotheses.'' The trilemma and fold theorems require caller-supplied hypotheses.
\item It does \textbf{not} claim that the refinement branch always produces a strictly finer regime---only that the structural \emph{form} of the response is one of the three lawful alternatives.
\end{itemize}

\subsection{Residual-driven improvement}

The theory does not identify all improvement with residual response. It does, however, isolate residual-driven refinement and restructuring as the deepest endogenous form of reflexive development presently formalized in this science. When a reflexive system refines its observational regime, shifts its regulatory architecture, or restructures its closure discipline because standing residual demands it, the improvement is driven by the system's own post-completion burden---not externally imposed. In this sense, the nonexhaustibility established earlier is not merely a limit but the source of lawful reflexive development.

\section{The general science of reflexive systems}

With both the static and dynamical arcs in view, we can state the field law.

\begin{keybox}[title=Master law of reflexive systems (canonical formulation)]
Strong internal completion in reflexive systems is structurally limited; its failure is lawful rather than catastrophic; and the resulting residual structure is organized, witnessable, enrichable, stratified by semantic barrier regime, and subject to lawful dynamical response---refinement, regime shift, or fold-obstruction transition---under standing residual burden.
\end{keybox}

\begin{theorem}[General law of reflexive systems]
Reflexive systems are governed by a common structural pattern:
\begin{enumerate}[label=(\roman*)]
\item strong internal completion is limited,
\item completion failure has lawful residual aftermath,
\item residual structure admits witness and enrichment forms,
\item semantic barrier regimes stratify reflexive closure,
\item these regimes can integrate back into the residual theory,
\item and dynamical response to standing residual is itself lawful: refinement, proper regime shift, or bookkeeping reconfiguration exhaust the structural alternatives, with fold obstruction forcing closure obstruction when iterate geometry fails, and adequacy burden relocating rather than vanishing.
\end{enumerate}
\end{theorem}

\subsection{What is unified here}

In one arc: completion modes \(\to\) anchored impossibility \(\to\) forced aftermath \(\to\) witnesses and enrichment \(\to\) semantic summit stratification \(\to\) integration \(\to\) lawful dynamical response. The \emph{static} half (completion through integration) says what reflexive systems cannot do and what structure is forced; the \emph{dynamical} half (the Reflexive Development Law) says what happens next under honest barrier discipline. Together they constitute a theory of reflexive systems that \emph{persist and develop}, not merely systems that fail once.

\subsection{Relation to earlier traditions}

Gödel, Tarski, Kleene, and Lawvere (among others) established that diagonal and reflexive structure carries unavoidable cost. This paper's contribution is not to discard those costs but to show how they organize a \emph{field}: completion modes, joint barriers, lawful residue, semantic summit ordering, explicit integration, and lawful dynamical response. The advance is ecological---a theory of how reflexive systems live \emph{under} limitation and continue to develop---not a single sharper paradox.

\subsection{Beyond this formal realization}

Nothing in the field law requires that reflexive systems be studied only in the present abstract architecture, only in one theorem prover, or only in one programmatic stack. What matters scientifically is the \textbf{recurring pattern}: reflexive systems press toward honest internal completion, are blocked by lawful barrier families, \textbf{persist through structured residue} whose semantic strength is regime-dependent, and \textbf{respond to that residue under lawful dynamical constraints}. Any future formalism that studies reflexive systems at comparable depth will still have to confront analogues of these structures, or else explain why its subject falls outside this scientific niche.

\section{Boundaries and open frontiers}

\subsection{What is not claimed}

The paper does not claim:
\begin{itemize}
\item that all reflexive success is anchored success;
\item that off-anchor or relocated success collapses to the same theorem subject;
\item that unrestricted intuitionistic four-way cover follows from bare `Genuine` alone under the current definition;
\item that all semantic summits collapse into one frame;
\item that all future taxonomy or integration questions are settled;
\item or that every residual forces nontrivial refinement without side hypotheses.
\end{itemize}

These boundaries are scientific strengths, not weaknesses. They prevent the theory from overspeaking.

\subsection{Open frontiers}

Natural next frontiers include:
\begin{itemize}
\item deeper engine-native payload transport,
\item stronger internalization of the triple-augment integration seam,
\item richer taxonomy comparison work,
\item additional semantic barrier regimes beyond the present computational/proof-theoretic pair,
\item unbounded dynamical iteration: iterated application of the response trilemma across multi-step reflexive histories,
\item and explicit formalization of \emph{residual-driven improvement} as a disciplined refinement operator (beyond the structural trilemma proved here).
\end{itemize}

\subsection{Why the science already matters}

A science need not close every future summit to become real. It must identify the right subject, the right laws, the right regimes, the right aftermath structure, and the right dynamical response constraints. This paper does that. The payoff for readers outside the present formal stack is exactly that pattern---in proof assistants, reflective AI architectures, internal certification pipelines, and epistemic institutions---not the proprietary naming of any single module.

\section{Objections and clarifications}

\subsection{``Is this just incompleteness in new words?''}

No. Classical incompleteness is largely a \emph{negative} lesson about what cannot be proved or defined. This paper retains that moral but adds a \emph{positive} field law: lawful aftermath, typed witnesses and enrichment, semantic summit stratification, and explicit summit-to-residue integration. A reader who collapses all of that to ``another undecidability story'' has truncated the subject.

\subsection{``Is anchored completion too strong?''}

It is the strongest notion that still answers the intuitive question ``did the system actually settle its reflexive burden where it claims to?'' Weaker targets are scientifically legitimate but different. The paper separates anchored from relocated success precisely so that strength cannot be mistaken for loophole.

\subsection{``Is residual structure just fancy incompleteness remainder?''}

Incompleteness traditionally emphasizes what cannot be proved or defined. The aftermath layer here insists that \emph{positive} structure is forced: profiles, witnesses, typed payloads, and enrichment. Residue is organized scientific data, not mere ``what's left when you give up.''

\subsection{``Are the semantic summits merely examples?''}

They are the two principal summit regimes currently developed with full barrier interfaces in this science. They are exemplary and structurally central; they are not asserted to exhaust all future semantic geometry (\S7.1).

\subsection{``Why should anyone outside your formal stack care?''}

Because the problem form---representation, generation, certification, and honest completion---is ubiquitous. The formal stack is one audit path. The same pattern appears in \textbf{interactive theorem proving} and large formal libraries, in \textbf{reflective AI} architectures that score or police their own outputs, in \textbf{semantic verification} and \textbf{internal certification} pipelines, and in \textbf{epistemic institutions} that codify their own methods and failure modes. The law is meant to be readable from definitions and theorem surfaces without consulting source code.

\subsection{``Is the dynamical layer just hand-waving about `what happens next'?''}

No. The Reflexive Development Law (Theorem~\ref{theorem:rdl-trilemma}) is a proved structural trilemma, not a narrative suggestion. The fold obstruction theorem (Theorem~\ref{theorem:fold-obstruction}) is a machine-checked consequence of iterate geometry. The adequacy burden bridge is a family of proved transport lemmas. The dynamical layer claims \emph{exactly} what its hypotheses support: that response to standing residual takes one of three lawful forms, and that iterate-backed closure regimes face explicit obstruction. It does not claim that every residual is generative or that self-improvement is automatic.

\subsection{``Does the theory depend on one formalization?''}

The theorems as stated are proved in a specific formal realization, but the philosophical and scientific divisions (anchored vs relocated, completion vs residual, predicated vs full summit, static vs dynamical) are meant to survive implementation change. Alternative foundations should either vindicate the same moral or explain precisely which hypothesis fails.

\section{Conclusion}

\noindent Under the three barriers, anchored completion is impossible, yet organized residue, semantic stratification, summit integration, and the Reflexive Development Law complete one scientific arc. The static half says what reflexive systems cannot do and what structure is forced; the dynamical half says what happens next under honest barrier discipline.

Reflexive systems are governed not by the fantasy of completed self-transparency, but by the lawful interplay of completion pressure, barriered failure, structured residuality, semantic stratification, and residual-driven development.

\appendix

\section{Theorem map}

\begin{enumerate}[label=\textbf{T\arabic*.}, leftmargin=2.2em]
\item Barriered anchored completion is impossible.
\item Completion failure forces positive residual aftermath.
\item Barrier-linked \(R_4\) witnesses exist at the abstract level.
\item Residual witnesses admit single-column and mixed enrichment.
\item The computational semantic route yields a predicated summit.
\item The proof-theoretic route yields a full semantic barrier summit.
\item The computational summit does not upgrade for free to the proof-theoretic summit.
\item The proof-theoretic summit integrates into the enriched residual line through constant, repr-dependent, and triple-augment seams.
\item (Reflexive Development Law) Under standing residual burden, response to residual takes one of three lawful forms: refinement (preserving typed residual in a new observational regime), proper regime shift (nontrivially distinct architecture), or bookkeeping reconfiguration (same underlying structure).
\item (Fold obstruction) When an iterate-backed closure regime cannot reach its target by any closure operator's iteration, closure obstruction is forced; naive iterate adequacy is not a free property of all architectures (sharp countermodel).
\item (Adequacy burden relocation) The forced admissible \(R_4\) aftermath constitutes standing residual burden that transports across architecture casts and regime-snapshot alignment, rather than vanishing under honest dynamical episodes.
\end{enumerate}

\section{Formal artifact note}

The formal realization is machine-checked in \Lean{} with \Mathlib{}, across the reflexive-architecture-nonexhaustibility and \texttt{nems-lean} stacks, with 0~\texttt{sorry} and 0~\texttt{axiom} in project-owned modules.

\noindent Theorem-to-artifact correspondences:
\begin{itemize}[leftmargin=1.5em, itemsep=2pt]
\item \textbf{T1} (anchored impossibility): \leanref{barriered\_architecture\_admits\_no\_true\_honest\_anchored\_internal\_completion} in \texttt{PaperDAnchoredChain.lean}
\item \textbf{T2} (structured aftermath): \leanref{positive\_residual\_profile\_of\_triple\_barriers\_at\_anchor} in \texttt{PostFailureResidual.lean}; \leanref{honest\_aftermath\_carries\_admissible\_r4} in \texttt{Adequacy.lean}
\item \textbf{T3} (barrier-linked witnesses): \leanref{d002\_barrier\_linked\_r4\_witness\_holds} in \texttt{D002ResidualWitnessTarget.lean}
\item \textbf{T4} (mixed enrichment): \leanref{mixedTriplePayloadPromotionBridge} in \texttt{FromMixedTriple.lean}
\item \textbf{T5} (predicated summit): \texttt{KleenePredicatedResidualSummit.lean}
\item \textbf{T6} (proof-theoretic summit): \texttt{GodelProvBic.lean} in \texttt{nems-lean}
\item \textbf{T7} (summit comparison): \texttt{SummitComparison.lean} in \texttt{nems-lean}
\item \textbf{T8} (O7 integration): \texttt{ProvBicU123EnrichedR4Alignment.lean} (T1--T3 seams)
\item \textbf{T9} (response trilemma): \leanref{reflexive\_development\_law\_standing\_trilemma} in \texttt{ResidualDynamics.lean}
\item \textbf{T10} (fold obstruction): \leanref{reflexive\_development\_law\_fold\_obstruction\_iterate\_geometry} in \texttt{ResidualDynamics.lean}
\item \textbf{T11} (burden relocation): \leanref{d3\_adequacy\_aftermath\_standing\_burden} in \texttt{Adequacy.lean}
\end{itemize}

\section{Layer distinctions}

The paper distinguishes:
\begin{enumerate}[label=(\alph*)]
\item anchored completion from relocated success,
\item completion modes from residual aftermath,
\item predicated semantic summits from full semantic summits,
\item frozen summit comparisons from future taxonomy ambitions,
\item the static arc (impossibility + aftermath + enrichment + semantic stratification + integration) from the dynamical arc (persistence + response trilemma + fold obstruction + adequacy burden relocation),
\item and proved bounded dynamics from unearned ``self-improvement'' slogans.
\end{enumerate}

These distinctions are not rhetorical. They are part of the subject matter itself.

\section*{References}

\noindent
Barwise, J., and Moss, L. \emph{Vicious Circles: On the Mathematics of Non-Wellfounded Phenomena}. CSLI Publications, 1996.

\medskip
Boolos, G. \emph{The Logic of Provability}. Cambridge University Press, 1993.

\medskip
Feferman, S. ``Arithmetization of metamathematics in a general setting.'' \emph{Fundamenta Mathematicae} 49 (1960): 35--92.

\medskip
Gödel, K. ``On formally undecidable propositions of \emph{Principia Mathematica} and related systems I.'' 1931. English translation in \emph{Kurt Gödel: Collected Works}, Vol. I.

\medskip
Hofstadter, D. R. \emph{Gödel, Escher, Bach: An Eternal Golden Braid}. Basic Books, 1979.

\medskip
Kleene, S. C. \emph{Introduction to Metamathematics}. North-Holland, 1952.

\medskip
Lawvere, F. W. ``Diagonal arguments and cartesian closed categories.'' In \emph{Category Theory, Homology Theory and their Applications II}, Lecture Notes in Mathematics 92, Springer, 1969.

\medskip
Orseau, L., and Ring, M. ``Self-modification and mortality in artificial agents.'' In \emph{Artificial General Intelligence}, Lecture Notes in Computer Science 6830, Springer, 2011.

\medskip
Rogers, H. \emph{Theory of Recursive Functions and Effective Computability}. MIT Press, 1987.

\medskip
Smullyan, R. \emph{Gödel's Incompleteness Theorems}. Oxford University Press, 1992.

\medskip
Spivack, N. \emph{The General Science of Reflexive Systems: formal artifact manifests and machine-checked integration stack}. Unpublished machine-checked research artifact, 2026.

\medskip
Sturm, T. ``Self-knowledge and self-models.'' In \emph{The Oxford Handbook of the Self}, Oxford University Press, 2011.

\medskip
Tarski, A. ``The concept of truth in formalized languages.'' In \emph{Logic, Semantics, Metamathematics}, 2nd ed., Hackett, 1983.

\medskip
Yanofsky, N. S. ``A universal approach to self-referential paradoxes, incompleteness and fixed points.'' \emph{Bulletin of Symbolic Logic} 9.3 (2003): 362--386.

\end{document}